\section{الگوریتم‌های یادگیری ماشین}
\subsection{جنگل تصادفی}
جنگل تصادفی، مجموعه‌ای از درخت‌های تصمیم است که با نمونه‌گیری بازگشتی از آموزش و انتخاب تصادفی زیرمجموعۀ ویژگی‌ها در گره‌ها، هم‌بستگی خطای درخت‌ها را کاهش می‌دهد. ریشۀ این روش در ترکیب نمونه‌گیری \en{bootstrap} و ایده‌های مجموعه‌مدل است و صورت شناخته‌شدۀ آن با کار بریمن گسترش یافت \cite{ref103,ref104,ref108}. هدف، کاهش واریانس درخت منفرد بدون الزام به ساخت یک مدل خطی است.

اگر درخت $b$ احتمال $p_b(y=1\mid x)$ را برآورد کند، احتمال خروجی جنگل با $B$ درخت به‌صورت میانگین آن‌ها محاسبه می‌شود:
\begin{equation}
p_{\mathrm{RF}}(x)=\frac1B\sum_{b=1}^{B}p_b(y=1\mid x),\qquad
\widehat y=\mathbf1\{p_{\mathrm{RF}}(x)\ge\tau\}.
\end{equation}
این احتمال لزوماً احتمال بالینیِ کالیبره‌شده نیست. تنظیم آستانه، برچسب خروجی را تغییر می‌دهد و به‌خودی‌خود کالیبراسیون احتمال انجام نمی‌دهد. در نسخۀ اجراشده، جنگل ۲۰۰ درخت دارد، بذر ۴۲ و تعداد پردازشگر مدل دو است؛ برای آن جست‌وجوی فراپارامتر مستقل تعریف نشده است.

مدل‌های درختی به مقیاس یکسان ویژگی‌ها وابستگی مدل‌های فاصله‌محور را ندارند و می‌توانند برهم‌کنش غیرخطی مختصات بردار سند را یاد بگیرند. در مقابل، تعداد زیاد درخت‌ها هزینۀ آموزش و حافظه را افزایش می‌دهد؛ شاخص اهمیت ویژگی نیز برای مختصات یک بازنمایی توزیعی، به‌سادگی به واژۀ توضیح‌پذیر تبدیل نمی‌شود. علاوه بر این، جنگل بر نمونۀ کوچک محلی ممکن است ساختار معنایی کم‌کیفیت \en{Word2Vec} را جبران نکند.
\begin{figure}[htbp]
\centering\begin{tikzpicture}[
  node distance=5mm,
  stage/.style={draw=teal!60!black,fill=teal!3,rounded corners=3pt,
    inner sep=6pt,font=\small},
  flow/.style={-{Latex[length=2.2mm]},thick,draw=teal!60!black}]
\node[stage] (data) {\diagramtext{بردارهای صدبعدی سند از \en{Word2Vec} آموخته‌شده بر آموزش}};
\node[stage,below=of data] (fit) {\diagramtext{برازش ۲۰۰ درخت با نمونه‌گیری بازگشتی\\
  در هر گره، جست‌وجوی تقسیم میان زیرمجموعۀ تصادفی ویژگی‌ها}};
\node[stage,below=of fit] (mean) {\diagramtext{احتمال کلاس مثبت برای ورودی جدید\\
  $p_{\mathrm{RF}}(x)=\frac1{200}\sum_{b=1}^{200}p_b(y=1\mid x)$}};
\node[stage,below=of mean] (threshold) {\diagramtext{انتخاب آستانه با $F_1$ بر اعتبارسنجی\\
  $\widehat y=\mathbf1\{p_{\mathrm{RF}}(x)\ge\tau\}$؛ تثبیت پیش از آزمون}};
\foreach \a/\b in {data/fit,fit/mean,mean/threshold}\draw[flow] (\a) -- (\b);
\end{tikzpicture}

\caption{جریان جنگل تصادفی و انتخاب آستانه بر اعتبارسنجی در آزمایش اصلاح‌شده}
\label{fig:rf-flow}
\end{figure}

\subsection{درخت تصمیم}
درخت تصمیم با تقسیم بازگشتی فضای ویژگی‌ها، مسیر تصمیم هر نمونه را از ریشه به برگ تعیین می‌کند. پژوهش‌های هانت، کوئینلن و چارچوب \en{CART} از مبانی تاریخی این خانواده هستند \cite{ref124,ref125,ref126,ref127,ref128}. در طبقه‌بندی دودویی، احتمال کلاس در برگ از نسبت نمونه‌های آن کلاس در برگ برآورد می‌شود. معیار ناخالصی جینی برای نسبت‌های $p_0,p_1$ برابر است با
\begin{equation}
G=1-p_0^2-p_1^2,
\end{equation}
و آنتروپی به‌صورت $H=-\sum_{c\in\{0,1\}}p_c\log p_c$ تعریف می‌شود. یک تقسیم، کاهش وزن‌دار ناخالصی فرزندان نسبت به والد را بیشینه می‌کند. عمق زیاد و برگ‌های کم‌نمونه می‌توانند به بیش‌برازش منجر شوند؛ محدود کردن عمق یا افزایش حداقل اندازۀ برگ از راه‌های کنترل آن است \cite{ref137,ref138}.

در نوت‌بوک اصلاح‌شده، معیار از میان \en{gini}، \en{entropy} و \en{log\_loss}، عمق از میان ۳، ۵، ۸، ۱۰، ۱۵، ۲۰ و بدون محدودیت، و حداقل تعداد نمونه برای تقسیم از میان ۲، ۵، ۱۰، ۲۰ و ۵۰ انتخاب می‌شود. حداقل اندازۀ برگ، زیرمجموعۀ ویژگی‌ها و وزن کلاس نیز جزو فضای جست‌وجو هستند. فقط ۱۰ ترکیب با بذر ثابت جست‌وجو می‌شوند؛ بنابراین این آزمایش، جست‌وجوی جامع تمام ترکیب‌ها نیست. پارامترهای انتخاب‌شده در خروجی اجرای واقعی ثبت می‌شوند.

درخت منفرد قابلیت نمایش مسیر تصمیم دارد، اما وقتی ورودی، مختصات بردار \en{Word2Vec} است، تفسیر هر آستانه به یک مفهوم زبانی مشخص نیاز به تحلیل جداگانه دارد. حساسیت درخت به تغییر داده نیز دلیل مناسبی برای گزارش اعتبارسنجی و کنترل اندازۀ برگ است.
\begin{figure}[htbp]
\centering\begin{tikzpicture}[
  node distance=5mm,
  stage/.style={draw=teal!60!black,fill=teal!3,rounded corners=3pt,
    inner sep=6pt,font=\small},
  flow/.style={-{Latex[length=2.2mm]},thick,draw=teal!60!black}]
\node[stage] (search) {\diagramtext{انتخاب فراپارامتر از ۱۰ ترکیب با اعتبارسنجی سه‌بخشی روی آموزش\\
  بازنمایی هر تقسیم فقط بر بخش آموزشی همان تقسیم یاد گرفته می‌شود}};
\node[stage,below=of search] (split) {\diagramtext{ساخت بازگشتی درخت با بیشینه‌سازی کاهش ناخالصی\\
  توقف بر اساس عمق، اندازۀ برگ و دیگر قیدهای مدل}};
\node[stage,below=of split] (leaf) {\diagramtext{ورودی جدید از ریشه و شرط گره‌ها به برگ می‌رسد\\
  $p(y=1\mid x)=N_{+,\mathrm{leaf}}/N_{\mathrm{leaf}}$}};
\node[stage,below=of leaf] (out) {\diagramtext{انتخاب آستانه روی اعتبارسنجی مستقل\\
  احتمال برگ و آستانۀ تثبیت‌شده، برچسب نهایی را تعیین می‌کنند}};
\foreach \a/\b in {search/split,split/leaf,leaf/out}\draw[flow] (\a) -- (\b);
\end{tikzpicture}

\caption{فلوچارت مفهومی درخت تصمیم}
\label{fig:dt-flow}
\end{figure}

\subsection{تقویت گرادیانی و \en{XGBoost}}
\en{XGBoost} یک پیاده‌سازی مقیاس‌پذیر از تقویت درختی است که با افزودن پیاپی درخت‌ها، تابع زیان منظم‌شده را کاهش می‌دهد \cite{ref159}. مدل در گام $t$ از رابطۀ
\begin{equation}
F_t(x)=F_{t-1}(x)+\eta f_t(x)
\end{equation}
پیروی می‌کند. برای طبقه‌بندی لجستیک، احتمال خروجی از تبدیل سیگموئید امتیاز به دست می‌آید. هدف کلی شامل زیان داده و جریمۀ پیچیدگی درخت‌هاست:
\begin{equation}
\mathcal L=\sum_{i=1}^{N}\ell(y_i,F(x_i))+\sum_t\Omega(f_t),\qquad
\Omega(f)=\gamma T+\frac\lambda2\sum_{j=1}^{T}\omega_j^2.
\end{equation}
در این رابطه، $T$ تعداد برگ‌ها، $\omega_j$ وزن برگ، $\gamma$ هزینۀ افزودن برگ و $\lambda$ ضریب منظم‌سازی است. با تقریب مرتبۀ دوم زیان و تعریف گرادیان $g_i$ و هسین $h_i$، اگر $I_j$ مجموعۀ نمونه‌های برگ $j$ باشد، در حالت جریمۀ $L_2$ و نبود جریمۀ $L_1$ داریم:
\begin{equation}
G_j=\sum_{i\in I_j}g_i,\qquad H_j=\sum_{i\in I_j}h_i,\qquad
\omega_j^*=-\frac{G_j}{H_j+\lambda}.
\end{equation}
سود تقسیم یک برگ به دو فرزند نیز به‌صورت زیر نوشته می‌شود:
\begin{equation}
\mathrm{Gain}=\frac12\left(\frac{G_L^2}{H_L+\lambda}+
\frac{G_R^2}{H_R+\lambda}-\frac{(G_L+G_R)^2}{H_L+H_R+\lambda}\right)-\gamma.
\end{equation}
این روابط، نقش ساختار درخت و منظم‌سازی را نشان می‌دهند؛ تضمینی برای برتری یک الگوریتم روی هر پیکره ایجاد نمی‌کنند.

در مقایسۀ متمرکز، عمق درخت از مجموعۀ ۳، ۵، ۷ و ۱۰، تعداد برآوردگر از ۵۰، ۱۰۰ و ۱۵۰ و نرخ یادگیری از ۰٫۰۵، ۰٫۱ و ۰٫۲ جست‌وجو می‌شود. تابع هدف \en{binary:logistic}، معیار داخلی \en{logloss} و روش درخت \en{hist} است. در سناریوی مشارکتی، پارامترها ثابت و مستقل از این جست‌وجوی متمرکز تعریف شده‌اند؛ این دو تنظیم باید از یکدیگر متمایز بمانند.
\begin{figure}[htbp]
\centering\begin{tikzpicture}[
  node distance=5mm,
  stage/.style={draw=teal!60!black,fill=teal!3,rounded corners=3pt,
    inner sep=6pt,font=\small},
  flow/.style={-{Latex[length=2.2mm]},thick,draw=teal!60!black}]
\node[stage] (start) {\diagramtext{آغاز از امتیاز پایۀ لجستیک $F_0$ و فراپارامترهای هر ترکیب}};
\node[stage,below=of start] (gradient) {\diagramtext{محاسبۀ گرادیان و هسین زیان برای نمونه‌های آموزش\\
  ساخت درخت با سود تقسیم و جریمۀ پیچیدگی}};
\node[stage,below=of gradient] (update) {\diagramtext{تعیین وزن برگ و افزودن درخت جدید\\
  $\omega_j=-G_j/(H_j+\lambda),\quad F_t=F_{t-1}+\eta f_t$}};
\node[stage,below=of update] (prob) {\diagramtext{پس از تعداد درخت تعیین‌شده، تبدیل امتیاز به احتمال\\
  $p(x)=[1+\exp(-F_T(x))]^{-1}$}};
\node[stage,below=of prob] (out) {\diagramtext{فراپارامترها از اعتبارسنجی متقابل آموزش انتخاب می‌شوند\\
  مدل و آستانه با اعتبارسنجی مستقل تثبیت می‌شوند}};
\foreach \a/\b in {start/gradient,gradient/update,update/prob,prob/out}\draw[flow] (\a) -- (\b);
\end{tikzpicture}

\caption{نمای مفهومی آموزش تقویتی \en{XGBoost}}
\label{fig:xgb-flow}
\end{figure}

\subsection{بیز ساده و رأی‌گیری احتمالی}
در بیز ساده، ویژگی‌ها مشروط به کلاس مستقل فرض می‌شوند. با استفاده از قضیۀ بیز، قاعدۀ طبقه‌بندی از رابطۀ
\begin{equation}
\widehat y=\arg\max_{c\in\{0,1\}}P(c)\prod_{j=1}^{d}P(x_j\mid c)
\end{equation}
به دست می‌آید. این استقلال معمولاً برای مختصات بازنمایی متن دقیقاً برقرار نیست، اما مدل همچنان می‌تواند یک خط مبنای محاسباتی مفید باشد \cite{ref112,ref142,ref176}. در \en{GaussianNB}، توزیع شرطی هر ویژگی با میانگین و واریانس کلاس مدل می‌شود:
\begin{equation}
P(x_j\mid c)=\frac{1}{\sqrt{2\pi\sigma_{cj}^{2}}}
\exp\!\left(-\frac{(x_j-\mu_{cj})^2}{2\sigma_{cj}^{2}}\right).
\end{equation}
در \en{BernoulliNB}، هر ویژگی به متغیر دودویی $z_j$ تبدیل می‌شود و سهم آن برابر $q_{cj}^{z_j}(1-q_{cj})^{1-z_j}$ است. هموارسازی لاپلاس از احتمال صفر برای رخدادهای مشاهده‌نشده جلوگیری می‌کند.

بردار میانگین \en{Word2Vec} ممکن است مؤلفه‌های منفی داشته باشد؛ بنابراین \en{MultinomialNB} مستقیماً مدل مناسبی برای همین ورودی نیست و در نسخۀ اصلاح‌شده استفاده نمی‌شود. ساختار اجراشده ابتدا \en{StandardScaler} را بر آموزش برازش می‌کند، سپس دو مدل گاوسی و برنولی را با رأی‌گیری نرم ترکیب می‌کند:
\begin{equation}
z_j=\frac{x_j-\mu_{j,\mathrm{train}}}{s_{j,\mathrm{train}}},\qquad
p_{\mathrm{vote}}(c\mid x)=\frac12\left(p_{\mathrm G}(c\mid z)+p_{\mathrm B}(c\mid z)\right).
\end{equation}
برای مدل برنولی، آستانۀ دودویی‌سازی صفر است و پارامتر هموارسازی یک است. \en{var\_smoothing} مدل گاوسی برابر $10^{-8}$ است. استانداردساز درون خط‌لوله و فقط روی آموزش همان تقسیم برازش می‌شود. آستانۀ تصمیم نهایی روی احتمال رأی‌گیری‌شده اعمال می‌شود و با آستانۀ دودویی‌سازی ویژگی تفاوت دارد.

نام این مجموعۀ دوتایی در گزارش \en{Naive Bayes voting} است تا صورت دقیق آزمایش روشن بماند.
\begin{figure}[htbp]
\centering\begin{tikzpicture}[
  node distance=5mm,
  stage/.style={draw=teal!60!black,fill=teal!3,rounded corners=3pt,
    inner sep=6pt,font=\small},
  flow/.style={-{Latex[length=2.2mm]},thick,draw=teal!60!black}]
\node[stage] (vector) {\diagramtext{بردار صدبعدی سند و استانداردسازی آموخته‌شده بر آموزش\\
  $z_j=(x_j-\mu_{j,\mathrm{train}})/s_{j,\mathrm{train}}$}};
\node[stage,below=of vector] (fit) {\diagramtext{برازش دو مدل روی همان بخش آموزش\\
  \en{GaussianNB}: واریانس هموارشده؛ \en{BernoulliNB}: دودویی‌سازی در صفر}};
\node[stage,below=of fit] (vote) {\diagramtext{ترکیب احتمال دو مدل با رأی‌گیری نرم و وزن برابر\\
  $p(c\mid x)=\tfrac12[p_{\mathrm G}(c\mid z)+p_{\mathrm B}(c\mid z)]$}};
\node[stage,below=of vote] (threshold) {\diagramtext{انتخاب آستانۀ احتمال روی اعتبارسنجی مستقل\\
  این آستانه با آستانۀ دودویی‌سازی ویژگی‌های برنولی متفاوت است}};
\foreach \a/\b in {vector/fit,fit/vote,vote/threshold}\draw[flow] (\a) -- (\b);
\end{tikzpicture}

\caption{ساختار رأی‌گیری نرم بیز گاوسی\lr{--}برنولی روی بازنمایی استانداردشدۀ سند}
\label{fig:bayes-model}
\end{figure}
